% Build: latex/build.sh labs/002-spark-graphframes/report/report.tex
\documentclass{iasa-lab}
\addbibresource{references.bib}

\labnumber{2}
\topic{Робота з графовими структурами}
\student{Гавлицький Іван}
\studentgroup{КІ-61мп}
\variant{3: Німеччина -- Велика Британія}

\begin{document}

\maketitle

\tableofcontents
\clearpage

% ============================================================================================
\anonsection{Мета роботи}

Ознайомитися з бібліотеками для обробки графів Spark GraphX та GraphFrames, вивчити основні алгоритми центральності та пошуку шляхів. Набути практичних навичок аналізу зв'язків у великих масивах даних, використовуючи набори даних авіаперевезень OpenFlights. Навчитися виконувати з'єднання даних (Join) для підготовки атрибутів вершин та ребер графа.

% ============================================================================================
\section{Опис задачі та вхідних даних}

\subsection{Постановка задачі}

Відповідно до завдання необхідно:
\begin{enumerate}
  \item завантажити дані OpenFlights у DataFrame, з'єднанням (Join) додати до маршрутів координати аеропортів, обчислити відстань кожного рейсу за формулою гаверсинуса та знайти авіакомпанії з найбільшою і найменшою сумарною довжиною рейсів;
  \item побудувати в Spark GraphFrames граф (вершини -- аеропорти, ребра -- рейси) і знайти всі маршрути між країнами варіанта з не більше ніж двома пересадками (до трьох сегментів), використавши пошук мотивів і BFS, та порівняти їх за результатами і складністю реалізації;
  \item знайти найкоротший і найдовший за фізичною відстанню (км) маршрут між двома аеропортами країн варіанта з не більше ніж трьома пересадками;
  \item знайти алгоритмом PageRank 5 найвпливовіших аеропортів, порівняти їх з найбільшими хабами світу та проаналізувати вплив параметра \texttt{resetProbability};
  \item виділити кластери аеропортів (Connected Components або Label Propagation), проаналізувати кластери з не менше ніж 5 аеропортів і знайти ізольовані підграфи з не менше ніж 2 вузлів.
\end{enumerate}

Варіант~3 сформульовано як «Німеччина -- Велика Британія». Маршрути в OpenFlights напрямлені, тому в завданні~2 його інтерпретовано як пошук маршрутів в одному напрямку: з аеропортів Німеччини до аеропортів Великої Британії. Зворотний напрямок не обчислювався. Для завдання~3 обрано пару аеропортів цих країн: Берлін-Тегель (TXL) -- Единбург (EDI).

\subsection{Вхідні дані}

Використано відкриті дані проєкту OpenFlights~\cite{openflights}. Файли завантажуються з репозиторію проєкту на GitHub, зафіксованого на одному коміті~\cite{openflights-github}, тому кожен запуск отримує ті самі дані. Структуру файлів наведено в табл.~\ref{tab:data}. Файли не мають рядка заголовка, відсутні значення позначено як \verb|\N|.

\begin{table}[H]
  \caption{Структура вхідних даних}\label{tab:data}
  \begin{tabularx}{\textwidth}{|l|r|Y|}
    \hline
    Файл & Записів & Основні поля \\
    \hline
    airports-extended.dat & 12\,668 & id, name, city, country, iata, icao, latitude, longitude, type \\
    airlines.dat & 6\,162 & id, name, iata, icao, country, active \\
    routes.dat & 67\,663 & airline, airline\_id, src, src\_id, dst, dst\_id, codeshare, stops, equipment \\
    \hline
  \end{tabularx}
\end{table}

За описом джерела~\cite{openflights}, база маршрутів відображає стан на червень 2014~року: 67\,663 маршрути між 3\,321 аеропортом. Маршрути напрямлені: якщо авіакомпанія літає з A до B і назад, це два записи. Поле \texttt{codeshare = 'Y'} позначає рейс, який виконує інший перевізник, а продає ця авіакомпанія під своїм кодом. Частоти рейсів, кількість місць і пасажиропотік у даних відсутні: маршрут лише вказує, що авіакомпанія сполучає два аеропорти. Файл \texttt{airports-extended.dat}, на відміну від \texttt{airports.dat}, містить також залізничні станції, поромні термінали та записи, додані користувачами.

% ============================================================================================
\section{Опис реалізації}

\subsection{Програмне середовище}

Роботу виконано в ноутбуці Jupyter (див. додаток~А) на кластері Apache Spark~4.0.4 у режимі Standalone~\cite{spark-docs}, розгорнутому локально в Docker Compose. Склад середовища наведено в табл.~\ref{tab:env}, схему взаємодії компонентів -- на рис.~\ref{fig:architecture}.

\begin{table}[H]
  \caption{Програмне середовище}\label{tab:env}
  \begin{tabularx}{\textwidth}{|l|Y|}
    \hline
    Компонент & Призначення та параметри \\
    \hline
    spark-master & менеджер ресурсів кластера, \texttt{spark://127.0.0.1:7077} \\
    spark-worker-1, -2 & виконавці (executors): по 4 ядра та 6 ГБ пам'яті, Python~3.13 для коду на Python \\
    Драйвер & процес ядра Jupyter: PySpark~4.0.4, GraphFrames~0.12.3, pandas \\
    Дані & директорія \texttt{data/}, змонтована за тим самим шляхом у кожен контейнер \\
    \hline
  \end{tabularx}
\end{table}

\addimg{../assets/architecture.pdf}{0.95}{Архітектура кластера}{fig:architecture}

Ноутбук не є клієнтом кластера, а виконує роль драйвера: PySpark запускає локальну JVM, яка будує план кожного запиту оптимізатором Catalyst і ділить його на задачі. Master лише виділяє ресурси (кроки 1--2 на рис.~\ref{fig:architecture}); executors у контейнерах підключаються до драйвера, самі читають дані з \texttt{data/} і обмінюються проміжними результатами (shuffle) (кроки 3--5). Контейнери використовують мережу хоста, тому executors звертаються до драйвера за адресою 127.0.0.1. Стан контейнерів після запуску показано на рис.~\ref{fig:cluster}, вебінтерфейс master -- на рис.~\ref{fig:spark-master}: два робочі вузли по 4 ядра, і всі 8 ядер зайняті застосунком \texttt{lab002} (сесією ноутбука).

\addshot{cluster-ps.png}{Контейнери кластера після запуску}{fig:cluster}

\addshot{spark-master.png}{Вебінтерфейс Spark Master: робочі вузли та запущений застосунок}{fig:spark-master}

\subsection{Завантаження та очищення даних}

Файли завантажує скрипт \texttt{download\_data.py} (додаток~Б, рис.~\ref{fig:download}). Оскільки файли не мають заголовків, схеми задано явно, а значення \verb|\N| читаються як \texttt{NULL} (лістинг~\ref{lst:read}). Шлях до файлів абсолютний: файл читають executors, а директорія \texttt{data/} змонтована за цим шляхом на кожному вузлі. Кількість записів показано на рис.~\ref{fig:row-counts}, приклад маршрутів -- на рис.~\ref{fig:routes-preview}.

\addshot{download.png}{Завантаження даних OpenFlights}{fig:download}

\noindent\begin{minipage}{\linewidth}
\begin{lstlisting}[language=Python, caption={Читання файлів OpenFlights зі схемами}, label={lst:read}]
S, I, D = StringType(), IntegerType(), DoubleType()


def schema(*fields):
    return StructType([StructField(name, dtype) for name, dtype in fields])


AIRPORTS = schema(
    ("id", I), ("name", S), ("city", S), ("country", S), ("iata", S), ("icao", S), ("latitude", D),
    ("longitude", D), ("altitude", I), ("timezone", D), ("dst", S), ("tz", S), ("type", S), ("source", S),
)
AIRLINES = schema(
    ("id", I), ("name", S), ("alias", S), ("iata", S), ("icao", S), ("callsign", S), ("country", S), ("active", S),
)
ROUTES = schema(
    ("airline", S), ("airline_id", I), ("src", S), ("src_id", I), ("dst", S), ("dst_id", I), ("codeshare", S),
    ("stops", I), ("equipment", S),
)


def read(name, file_schema):
    # Absolute path: executors read the file themselves, and data/ is mounted at this path on every worker.
    return spark.read.csv(str(OPENFLIGHTS / name), schema=file_schema, nullValue="\\N")


airports = read("airports-extended.dat", AIRPORTS)
airlines = read("airlines.dat", AIRLINES)
routes = read("routes.dat", ROUTES)
for name, df in [("airports", airports), ("airlines", airlines), ("routes", routes)]:
    df.createOrReplaceTempView(name)
\end{lstlisting}
\end{minipage}
\medskip

\addshot{row-counts.png}{Кількість записів у таблицях}{fig:row-counts}

\addshot{routes-preview.png}{Перші маршрути таблиці routes}{fig:routes-preview}

Перевірку якості маршрутів наведено в лістингу~\ref{lst:quality} і на рис.~\ref{fig:quality}: 423 маршрути не мають ідентифікатора аеропорту (з'єднання з координатами для них неможливе), 479 -- ідентифікатора авіакомпанії, 14\,597 є codeshare-маршрутами, а один маршрут є петлею (аеропорт вильоту збігається з аеропортом прильоту).

\noindent\begin{minipage}{\linewidth}
\begin{lstlisting}[language=SQL, caption={Перевірка якості маршрутів}, label={lst:quality}]
SELECT
    COUNT(*) AS routes,
    COUNT_IF(src_id IS NULL OR dst_id IS NULL) AS without_airport_id,
    COUNT_IF(airline_id IS NULL) AS without_airline_id,
    COUNT_IF(codeshare = 'Y') AS codeshare,
    COUNT_IF(src_id = dst_id) AS same_airport

FROM
    routes
\end{lstlisting}
\end{minipage}
\medskip

\addshot{route-quality.png}{Проблеми якості в таблиці маршрутів}{fig:quality}

Відстань по великому колу між аеропортами вильоту $(\varphi_1, \lambda_1)$ і прильоту $(\varphi_2, \lambda_2)$ обчислюється за формулою гаверсинуса:
\begin{equation}\label{eq:haversine}
  d = 2R \arcsin\sqrt{\sin^2\frac{\varphi_2 - \varphi_1}{2} + \cos\varphi_1 \cos\varphi_2 \sin^2\frac{\lambda_2 - \lambda_1}{2}}, \quad R = 6371~\text{км}.
\end{equation}

Представлення \texttt{route\_distances\_all} (лістинг~\ref{lst:distances}) з'єднує маршрути з таблицею аеропортів двічі -- для аеропорту вильоту та прильоту -- і обчислює відстань~(\ref{eq:haversine}). \texttt{INNER JOIN} відкидає маршрути без ідентифікатора аеропорту, умова \texttt{WHERE} -- петлю.

\noindent\begin{minipage}{\linewidth}
\begin{lstlisting}[language=SQL, caption={Координати маршрутів і відстань за формулою гаверсинуса}, label={lst:distances}]
CREATE OR REPLACE TEMP VIEW route_distances_all AS
SELECT
    r.airline,
    r.airline_id,
    r.codeshare,
    r.src_id,
    r.dst_id,
    s.iata AS src_iata,
    d.iata AS dst_iata,
    s.country AS src_country,
    d.country AS dst_country,
    2 * 6371 * ASIN(SQRT(
        POWER(SIN(RADIANS(d.latitude - s.latitude) / 2), 2)
        + COS(RADIANS(s.latitude)) * COS(RADIANS(d.latitude))
        * POWER(SIN(RADIANS(d.longitude - s.longitude) / 2), 2)
    )) AS distance_km

FROM
    routes AS r

INNER JOIN
    airports AS s
    ON r.src_id = s.id

INNER JOIN
    airports AS d
    ON r.dst_id = d.id

WHERE 1=1
    AND r.src_id != r.dst_id
\end{lstlisting}
\end{minipage}
\medskip

Сортування за відстанню виявило пошкоджений запис (рис.~\ref{fig:longest-raw},~\ref{fig:broken-airport}): найдовші «рейси» (понад 15\,900~км, довше за будь-який реальний безпосадковий рейс) ведуть до аеропорту 5613, у якого місто й країна -- Солвезі, Замбія, а назва, код ICAO та координати належать аеродрому Los Alamitos у Каліфорнії. Маршрути через цей аеропорт виключено (лістинг~\ref{lst:clean}); наступний за довжиною маршрут SYD $\to$ DFW (13\,808~км) реальний.

\addshot{longest-raw.png}{Найдовші маршрути до очищення}{fig:longest-raw}

\addshot{broken-airport.png}{Пошкоджений запис аеропорту 5613}{fig:broken-airport}

\noindent\begin{minipage}{\linewidth}
\begin{lstlisting}[language=SQL, caption={Виключення маршрутів через пошкоджений аеропорт}, label={lst:clean}]
CREATE OR REPLACE TEMP VIEW route_distances AS
SELECT
    *

FROM
    route_distances_all

WHERE 1=1
    AND src_id NOT IN (5613)
    AND dst_id NOT IN (5613)
\end{lstlisting}
\end{minipage}
\medskip

Після очищення залишилося 67\,235 маршрутів з координатами (рис.~\ref{fig:distance-stats}) довжиною від 1 до 13\,808~км, у середньому 1\,845~км.

\addshot{distance-stats.png}{Статистика відстаней після очищення}{fig:distance-stats}

\subsection{Обґрунтування вибору алгоритмів}

\textbf{Модель графа.} Вершини -- аеропорти, що мають хоча б один маршрут, ребра -- напрямки «аеропорт $\to$ аеропорт» з атрибутами \texttt{distance\_km} і кількістю авіакомпаній. Маршрути різних авіакомпаній між тією ж парою аеропортів об'єднано в одне ребро: інакше кожен шлях повторювався б для кожної комбінації перевізників, а пошук маршрутів вимірює зв'язність мережі, а не кількість перевізників. Граф напрямлений, як і маршрути OpenFlights.

\textbf{Пошук мотивів (завдання 2, 3).} Завдання вимагають \emph{усіх} маршрутів до заданої кількості пересадок і їх фізичної довжини. Мотив \texttt{(a)-[e1]->(b); (b)-[e2]->(c)}~\cite{graphframes-docs} перетворюється на з'єднання таблиць ребер і вершин, тому дає кожен шлях заданої довжини з атрибутами всіх ребер, з яких сумується відстань. Для мотивів написано допоміжну функцію \texttt{simple\_paths} (додаток~Б): вона будує ланцюжок з $k$ ребер, додає умови на першу й останню вершину та забороняє повторне відвідування аеропорту (простий шлях).

\textbf{BFS (завдання 2).} Пошук у ширину -- рекомендований завданням альтернативний метод; його виконано для порівняння з мотивами. Він знаходить шляхи з мінімальною кількістю сегментів, але не всі шляхи до заданої довжини і без урахування кілометрів, тому для завдання~3 не підходить.

\textbf{PageRank (завдання 4)}~\cite{needham2019} оцінює транзитивний вплив: ранг аеропорту залежить не лише від кількості напрямків, а й від рангу аеропортів, з яких до нього летять. Для порівняння обчислено ступеневу центральність (кількість напрямків).

\textbf{Connected Components і Label Propagation (завдання 5).} Слабко зв'язні компоненти (напрямок ребер ігнорується) прямо відповідають на питання про ізольовані підграфи: група аеропортів, не досяжна з основної мережі, -- окрема компонента. Проте майже всі аеропорти утворюють одну компоненту, тому для виділення кластерів усередині мережі додатково застосовано Label Propagation: кожна вершина ітеративно переймає найчастішу мітку сусідів, і тісніше пов'язані між собою аеропорти отримують спільну мітку.

\subsection{Завдання 1. Сумарна довжина рейсів авіакомпаній}

Сумарну довжину рейсів обчислює запит з лістингу~\ref{lst:airlines}. Враховано лише маршрути, які авіакомпанія виконує сама: codeshare-маршрут -- це рейс іншого перевізника. Кожен маршрут рахується один раз, незалежно від частоти рейсів. Результати для порядку \texttt{DESC} і \texttt{ASC} наведено на рис.~\ref{fig:airlines-top},~\ref{fig:airlines-bottom}.

\noindent\begin{minipage}{\linewidth}
\begin{lstlisting}[language=SQL, caption={Сумарна довжина маршрутів авіакомпаній}, label={lst:airlines}]
SELECT
    a.name AS airline,
    a.country,
    COUNT(*) AS routes,
    ROUND(SUM(f.distance_km)) AS total_km

FROM
    route_distances AS f

INNER JOIN
    airlines AS a
    ON f.airline_id = a.id

WHERE 1=1
    AND f.codeshare IS NULL

GROUP BY a.id, a.name, a.country

ORDER BY total_km DESC

LIMIT 10
\end{lstlisting}
\end{minipage}
\medskip

\addshot{airlines-top.png}{Авіакомпанії з найбільшою сумарною довжиною маршрутів}{fig:airlines-top}

\addshot{airlines-bottom.png}{Авіакомпанії з найменшою сумарною довжиною маршрутів}{fig:airlines-bottom}

\subsection{Завдання 2. Пошук маршрутів Німеччина -- Велика Британія}

Граф будується з представлення \texttt{route\_distances} (лістинг~\ref{lst:graph}) і містить 3\,329 вершин і 37\,269 ребер (рис.~\ref{fig:graph}). Обидва DataFrame кешуються, бо кожен наступний алгоритм читає їх багато разів.

\begin{lstlisting}[language=Python, caption={Побудова графа авіасполучень}, label={lst:graph}]
edges = spark.sql(r"""
SELECT
    src_id AS src,
    dst_id AS dst,
    FIRST(distance_km) AS distance_km,
    COUNT(DISTINCT airline_id) AS airlines

FROM
    route_distances

GROUP BY src_id, dst_id
""").cache()
edges.createOrReplaceTempView("edges")

vertices = spark.sql(r"""
SELECT
    a.id,
    COALESCE(a.iata, a.icao, CAST(a.id AS STRING)) AS code,
    a.name,
    a.city,
    a.country,
    a.latitude,
    a.longitude

FROM
    airports AS a

WHERE 1=1
    AND a.id IN (SELECT src FROM edges UNION SELECT dst FROM edges)
""").cache()
vertices.createOrReplaceTempView("vertices")

g = GraphFrame(vertices, edges)
\end{lstlisting}

\addshot{graph.png}{Розмір графа авіасполучень}{fig:graph}

\subsubsection{Пошук мотивів}

Маршрути з 1, 2 і 3 сегментами знаходить функція \texttt{simple\_paths} (лістинг~\ref{lst:motifs}). Для трьох сегментів вона формує мотив \texttt{(a)-[e1]->(b); (b)-[e2]->(c); (c)-[e3]->(d)} і умову \texttt{a.country = 'Germany' AND d.country = 'United Kingdom'} разом із шістьма нерівностями \texttt{a.id != b.id}, \texttt{a.id != c.id}, \ldots{} для всіх пар вершин. Кількість знайдених маршрутів наведено на рис.~\ref{fig:motif-counts}, найкоротші маршрути кожної довжини -- на рис.~\ref{fig:motif-examples}.

\noindent\begin{minipage}{\linewidth}
\begin{lstlisting}[language=Python, caption={Пошук маршрутів мотивами}, label={lst:motifs}]
FROM_DE = "{v}.country = 'Germany'"
TO_UK = "{v}.country = 'United Kingdom'"

de_uk = {segments: simple_paths(g, segments, FROM_DE, TO_UK).cache() for segments in (1, 2, 3)}
\end{lstlisting}
\end{minipage}
\medskip

\addshot{motif-counts.png}{Кількість маршрутів з Німеччини до Великої Британії за кількістю сегментів}{fig:motif-counts}

\addshot{motif-examples.png}{Найкоротші маршрути з 1, 2 і 3 сегментами}{fig:motif-examples}

Як Spark виконує мотив, видно з плану запиту у вебінтерфейсі драйвера (рис.~\ref{fig:spark-motif-dag}, мотив з одного сегмента). Таблиця ребер (37\,269 рядків) з'єднується з вершинами двічі: для аеропорту вильоту та прильоту. Оптимізатор Catalyst переніс умови на країну до з'єднань: з 3\,329 вершин фільтр залишає 32 аеропорти Німеччини та 52 Великої Британії. Ці невеликі таблиці розсилаються на всі executors (\texttt{BroadcastExchange}), тож з'єднання відбувається без перерозподілу (shuffle) таблиці ребер: перше з'єднання залишає 1\,409 ребер з Німеччини, друге -- 73 прямі маршрути.

\addshot{spark-motif-dag.png}{Фрагмент плану виконання мотиву з одного сегмента у вебінтерфейсі Spark}{fig:spark-motif-dag}

Аеропорти пересадки на маршрутах з однією пересадкою наведено на рис.~\ref{fig:transfer-hubs}.

\addshot{transfer-hubs.png}{Найчастіші аеропорти пересадки (2 сегменти)}{fig:transfer-hubs}

\subsubsection{Пошук у ширину (BFS)}

Той самий пошук методом BFS виконує один виклик (лістинг~\ref{lst:bfs}). BFS повертає всі шляхи мінімальної довжини (за кількістю сегментів), але мінімум визначається один для всієї множини початкових вершин: пошук іде одночасно від усіх аеропортів Німеччини і зупиняється на першій глибині, на якій досягнуто будь-якого аеропорту Великої Британії. З частини німецьких аеропортів є прямі рейси, тому ця глибина дорівнює 1, і результат -- 73 прямі маршрути (рис.~\ref{fig:bfs}), ті самі, що й у мотиву з одного сегмента.

\noindent\begin{minipage}{\linewidth}
\begin{lstlisting}[language=Python, caption={Пошук маршрутів методом BFS}, label={lst:bfs}]
bfs = g.bfs(fromExpr="country = 'Germany'", toExpr="country = 'United Kingdom'", maxPathLength=3)
bfs.count(), bfs.columns
\end{lstlisting}
\end{minipage}
\medskip

\addshot{bfs.png}{Результат BFS для Німеччина $\to$ Велика Британія}{fig:bfs}

Отже, пошук від усієї Німеччини не показує, як BFS іде глибше, і не відрізняється від мотиву з одного сегмента. Тому BFS запущено ще й від одного аеропорту, з якого немає прямих рейсів до Великої Британії. Запит обирає такий аеропорт з найбільшою кількістю напрямків, щоб маршрутів з пересадкою було якомога більше: це Фрідріхсгафен (FDH). Для нього мінімальна довжина -- 2 сегменти, і BFS повертає 90 маршрутів з однією пересадкою (рис.~\ref{fig:bfs-one}). Цих маршрутів немає в результаті пошуку від усієї Німеччини, хоча FDH -- німецький аеропорт. Порівняння кількості шляхів, знайдених обома методами, наведено на рис.~\ref{fig:bfs-vs-motifs}; аналіз -- у підрозділі~\ref{sec:motifs-vs-bfs}.

\addshot{bfs-one.png}{Найкоротші за відстанню з маршрутів BFS для FDH}{fig:bfs-one}

\addshot{bfs-vs-motifs.png}{Кількість шляхів: BFS і мотиви з 1--3 сегментів}{fig:bfs-vs-motifs}

\subsection{Завдання 3. Найкоротший та найдовший маршрути TXL -- EDI}

Між Берлін-Тегель (TXL) і Единбургом (EDI) шукаються прості шляхи з 1--4 сегментів (до 3 пересадок) тією ж функцією \texttt{simple\_paths} (лістинг~\ref{lst:txl-edi}); довжина маршруту -- сума відстаней сегментів у кілометрах. Прямого рейсу немає; маршрутів з 2, 3 і 4 сегментів знайдено 40, 2\,745 і 180\,069 (рис.~\ref{fig:txl-edi-counts}). Найкоротші та найдовші маршрути наведено на рис.~\ref{fig:txl-edi-extremes}.

\noindent\begin{minipage}{\linewidth}
\begin{lstlisting}[language=Python, caption={Маршрути TXL $\to$ EDI з 1--4 сегментів}, label={lst:txl-edi}]
SRC, DST = "TXL", "EDI"

paths = None
for segments in (1, 2, 3, 4):
    found = simple_paths(g, segments, f"{{v}}.code = '{SRC}'", f"{{v}}.code = '{DST}'")
    paths = found if paths is None else paths.unionByName(found)
paths = paths.cache()

show(paths.groupBy("segments").count().orderBy("segments"))
\end{lstlisting}
\end{minipage}
\medskip

\addshot{txl-edi-counts.png}{Кількість маршрутів TXL $\to$ EDI за кількістю сегментів}{fig:txl-edi-counts}

\addshot{txl-edi-extremes.png}{Найкоротші та найдовші маршрути TXL $\to$ EDI}{fig:txl-edi-extremes}

Найкоротший маршрут -- TXL $\to$ AMS $\to$ EDI (1\,245~км) з пересадкою в Амстердамі. Найдовший -- TXL $\to$ ORD $\to$ HKG $\to$ EWR $\to$ EDI (37\,808~км, 94\,\% довжини екватора): формально він задовольняє умови (не більше трьох пересадок, жоден аеропорт не повторюється), але практичного сенсу не має. Найдовші маршрути складаються з трансатлантичних і транстихоокеанських сегментів між найбільшими хабами, бо саме з них є найдовші безпосадкові рейси.

\subsection{Завдання 4. PageRank}

PageRank запущено з \texttt{resetProbability = 0.15}, що відповідає коефіцієнту демпфування 0,85, і 20 ітераціями (лістинг~\ref{lst:pagerank}). Для порівняння до результату додано вхідний і вихідний ступені вершин (рис.~\ref{fig:pagerank}); аеропорти з найбільшою кількістю напрямків показано на рис.~\ref{fig:degree}.

\noindent\begin{minipage}{\linewidth}
\begin{lstlisting}[language=Python, caption={PageRank і ступенева центральність}, label={lst:pagerank}]
ranks = g.pageRank(resetProbability=0.15, maxIter=20).vertices
degrees = g.inDegrees.join(g.outDegrees, "id")

show(
    ranks.join(degrees, "id")
    .select("code", "name", "city", "country", F.round("pagerank", 2).alias("pagerank"), "inDegree", "outDegree")
    .orderBy(F.desc("pagerank")),
    10,
)
\end{lstlisting}
\end{minipage}
\medskip

\addshot{pagerank.png}{10 аеропортів з найбільшим PageRank}{fig:pagerank}

\addshot{degree.png}{10 аеропортів з найбільшою кількістю напрямків вильоту}{fig:degree}

Код \texttt{ISL} належить стамбульському аеропорту Ататюрк: у 2019~році код \texttt{IST} перейшов до нового аеропорту Стамбула, і довідник аеропортів OpenFlights це відображає, тоді як маршрути зібрано раніше (2014~рік).

Вплив параметра \texttt{resetProbability} досліджено для значень 0,05--0,85 (лістинг~\ref{lst:reset}, рис.~\ref{fig:reset}).

\noindent\begin{minipage}{\linewidth}
\begin{lstlisting}[language=Python, caption={PageRank для різних значень resetProbability}, label={lst:reset}]
top_by_reset = {}
for reset in (0.05, 0.15, 0.3, 0.5, 0.85):
    top = g.pageRank(resetProbability=reset, maxIter=20).vertices.orderBy(F.desc("pagerank")).limit(5)
    top_by_reset[reset] = [f"{row.code} ({row.pagerank:.1f})" for row in top.collect()]
\end{lstlisting}
\end{minipage}
\medskip

\addshot{reset-probability.png}{5 найвпливовіших аеропортів для різних значень resetProbability}{fig:reset}

У вебінтерфейсі драйвера видно, як GraphFrames виконує алгоритми (рис.~\ref{fig:spark-sql-list}). Для PageRank граф перетворюється на RDD (\texttt{rdd at GraphFrame.scala}), ітерації виконує GraphX, а результат повертається у DataFrame (\texttt{count at PageRank.scala}); кожен з п'яти запусків з лістингу~\ref{lst:reset} дає дві такі операції. Компоненти зв'язності (\texttt{TwoPhase.scala}, наступний підрозділ) реалізовано серією запитів DataFrame із записом проміжних результатів у checkpoint-директорію (\texttt{parquet at TwoPhase.scala}).

\addshot{spark-sql-list.png}{Запити PageRank і Connected Components у вкладці SQL / DataFrame вебінтерфейсу Spark}{fig:spark-sql-list}

\subsection{Завдання 5. Кластеризація та аналіз зв'язності}

Слабко зв'язні компоненти обчислює метод \texttt{connectedComponents} (лістинг~\ref{lst:components}). Алгоритм потребує checkpoint-директорії, спільної для драйвера та всіх executors; її розміщено в змонтованій директорії \texttt{data/}. Розміри всіх компонент наведено на рис.~\ref{fig:components}, кластери з не менше ніж 5 аеропортів -- на рис.~\ref{fig:components-5}.

\noindent\begin{minipage}{\linewidth}
\begin{lstlisting}[language=Python, caption={Компоненти зв'язності}, label={lst:components}]
components = g.connectedComponents().cache()
component_sizes = components.groupBy("component").count().orderBy(F.desc("count"))
show(component_sizes)

show(component_sizes.filter("count >= 5"))
\end{lstlisting}
\end{minipage}
\medskip

\addshot{components.png}{Розміри компонент зв'язності}{fig:components}

\addshot{components-5.png}{Компоненти з не менше ніж 5 аеропортів}{fig:components-5}

Ізольовані підграфи -- компоненти з 2 і більше аеропортів, крім найбільшої (основної мережі) -- знаходить запит з лістингу~\ref{lst:isolated} (рис.~\ref{fig:isolated}).

\noindent\begin{minipage}{\linewidth}
\begin{lstlisting}[language=Python, caption={Ізольовані підграфи}, label={lst:isolated}]
main_component = component_sizes.first().component
isolated = (
    components.join(component_sizes, "component")
    .filter((F.col("component") != main_component) & (F.col("count") >= 2))
    .groupBy("component", "count")
    .agg(
        F.concat_ws(", ", F.sort_array(F.collect_list(F.concat_ws(" ", "code", F.concat(F.lit("("), "city", F.lit(")")))))).alias("airports"),
        F.concat_ws(", ", F.array_distinct(F.collect_list("country"))).alias("countries"),
    )
    .orderBy(F.desc("count"))
)
\end{lstlisting}
\end{minipage}
\medskip

\addshot{isolated.png}{Ізольовані підграфи}{fig:isolated}

Для виділення кластерів усередині основної мережі запущено Label Propagation з 5 ітераціями (лістинг~\ref{lst:lpa}); він знайшов 71 спільноту з не менше ніж 5 аеропортів (рис.~\ref{fig:lpa-count}). Для кожної спільноти обчислено розмір, кількість країн, головну країну з її часткою та п'ять аеропортів з найбільшою кількістю напрямків (рис.~\ref{fig:lpa-summary}).

\noindent\begin{minipage}{\linewidth}
\begin{lstlisting}[language=Python, caption={Label Propagation}, label={lst:lpa}]
communities = g.labelPropagation(maxIter=5).cache()
community_sizes = communities.groupBy("label").count()
community_sizes.filter("count >= 5").count()
\end{lstlisting}
\end{minipage}
\medskip

\addshot{lpa-count.png}{Кількість спільнот з не менше ніж 5 аеропортів}{fig:lpa-count}

\addshot{lpa-summary.png}{15 найбільших спільнот Label Propagation}{fig:lpa-summary}

% ============================================================================================
\section{Аналіз результатів}

\subsection{Авіакомпанія з найбільшим кілометражем}

П'ять авіакомпаній з найбільшою сумарною довжиною власних маршрутів наведено в табл.~\ref{tab:airlines}; середню довжину маршруту обчислено як відношення сумарної довжини до кількості маршрутів.

\begin{table}[H]
  \caption{Авіакомпанії з найбільшою сумарною довжиною маршрутів}\label{tab:airlines}
  \begin{tabularx}{\textwidth}{|Y|l|r|r|r|}
    \hline
    Авіакомпанія & Країна & Маршрутів & Сумарно, км & Середня, км \\
    \hline
    Ryanair & Ірландія & 2\,484 & 3\,700\,425 & 1\,490 \\
    US Airways & США & 1\,446 & 3\,477\,460 & 2\,405 \\
    American Airlines & США & 1\,265 & 3\,342\,406 & 2\,642 \\
    United Airlines & США & 905 & 3\,046\,248 & 3\,366 \\
    Delta Air Lines & США & 1\,146 & 3\,013\,769 & 2\,630 \\
    \hline
  \end{tabularx}
\end{table}

Найбільший кілометраж має Ryanair: 3,70~млн~км на 2\,484 маршрутах. Лідерство забезпечує не довжина рейсів (у середньому 1\,490~км -- найменше серед п'ятірки), а кількість маршрутів: лоукостер будує мережу «з точки в точку» між сотнями європейських аеропортів, і майже кожна пара його баз має окремий маршрут. Американські мережеві перевізники мають у 1,7--2,7 раза менше маршрутів, але вони довші (2\,405--3\,366~км) завдяки трансконтинентальним і міжнародним рейсам.

Найменший кілометраж має TransHolding System (Бразилія): 36~км на 2 маршрутах. Усі 10 авіакомпаній з найменшим кілометражем мають рівно 2 власні маршрути загальною довжиною 36--237~км; серед них вертолітні (Papillon Grand Canyon Helicopters, Helijet) та місцеві перевізники.

Метрика рахує кожен маршрут один раз, тому показує довжину мережі маршрутів, а не реальний наліт: частоти рейсів у OpenFlights відсутні. Codeshare-маршрути (14\,597) не враховано, інакше авіакомпанії отримували б кілометри рейсів партнерів.

\subsection{Порівняння пошуку мотивів і BFS}\label{sec:motifs-vs-bfs}

Результати обох методів наведено в табл.~\ref{tab:motifs-bfs} (для FDH у ноутбуці обчислено лише загальну кількість маршрутів мотивами), порівняння реалізації -- у табл.~\ref{tab:motifs-bfs-impl}.

\begin{table}[H]
  \caption{Кількість знайдених маршрутів}\label{tab:motifs-bfs}
  \begin{tabularx}{\textwidth}{|Y|r|r|r|r|r|}
    \hline
    \multirow{2}{=}{Пошук} & \multicolumn{4}{c|}{Мотиви} & \multirow{2}{*}{BFS} \\
    \cline{2-5}
     & 1 сегм. & 2 сегм. & 3 сегм. & Разом & \\
    \hline
    Німеччина $\to$ Велика Британія & 73 & 9\,178 & 506\,498 & 515\,749 & 73 \\
    FDH $\to$ Велика Британія & 0 & \multicolumn{2}{c|}{--} & 6\,028 & 90 \\
    \hline
  \end{tabularx}
\end{table}

\begin{table}[!htb]
  \caption{Порівняння реалізації пошуку мотивів і BFS}\label{tab:motifs-bfs-impl}
  \begin{tabularx}{\textwidth}{|l|Y|Y|}
    \hline
    Критерій & Мотиви & BFS \\
    \hline
    Що знаходить & усі шляхи заданої довжини & усі шляхи мінімальної кількості сегментів \\
    Виклики & окремий мотив для кожної довжини (1, 2, 3) & один виклик з \texttt{maxPathLength=3} \\
    Прості шляхи & явні умови $u \ne v$ для кожної пари вершин: 1, 3, 6 умов & гарантовано: найкоротший шлях не повторює вершин \\
    Кінцеві вершини & умови на першу та останню вершину & \texttt{fromExpr}, \texttt{toExpr} \\
    Вартість & $2k$ з'єднань для $k$ сегментів; результат зростає комбінаторно & з'єднання по одному рівню, зупинка на першій глибині з результатом \\
    Відстань маршруту & сума атрибутів ребер \texttt{e1}..\texttt{ek} & можлива, але лише для найкоротших шляхів \\
    \hline
  \end{tabularx}
\end{table}

\textbf{Результати.} Мотиви знаходять 73 прямі маршрути, 9\,178 з однією пересадкою і 506\,498 з двома. Кожен додатковий сегмент збільшує кількість маршрутів у 55--126 разів: у графі в середньому 11,2 напрямку з аеропорту, а через хаби з сотнями напрямків число продовжень зростає ще швидше. BFS повертає всі маршрути з мінімальною кількістю сегментів, причому мінімум один для всієї множини початкових аеропортів. Від усієї Німеччини це 1 сегмент, тому BFS дає ті самі 73 прямі маршрути, що й мотив з одного сегмента; маршрути з пересадкою до результату не потрапляють навіть для аеропортів без прямих рейсів, таких як FDH. Щоб порівняти методи на однаковій множині шляхів, обидва запущено від одного аеропорту FDH: його мінімум -- 2 сегменти, і BFS знаходить 90 маршрутів з однією пересадкою, тоді як мотиви з 1--3 сегментів -- 6\,028 маршрутів з однією або двома пересадками (прямих рейсів з FDH немає). Отже, BFS відповідає на питання «як дістатися з мінімальною кількістю пересадок», а всі маршрути до двох пересадок, яких вимагає завдання, дають лише мотиви.

\textbf{Складність реалізації.} BFS -- один виклик бібліотеки з умовами на початкові й кінцеві вершини. Мотиви потребують окремого шаблону для кожної довжини шляху і явних фільтрів: GraphFrames не забороняє повторювати вершину в мотиві, тому без умов $a \ne c$ та інших знайшлися б маршрути на кшталт DUS $\to$ AMS $\to$ DUS $\to$ LHR. Функція \texttt{simple\_paths} (додаток~Б) генерує ці шаблони та умови, але логіка «простого шляху» -- відповідальність користувача, а не бібліотеки. Обчислювальна вартість також різна: мотив з $k$ сегментів -- це $2k$ з'єднань (рис.~\ref{fig:spark-motif-dag}), і результат для трьох сегментів (506\,498 рядків) необхідно повністю матеріалізувати, тоді як BFS зупиняється, щойно знайдено ціль.

\textbf{Пересадки.} Найчастіші аеропорти пересадки -- курортні аеропорти Іспанії та Португалії: Пальма-де-Мальорка (648 маршрутів), Малага (391), Тенерифе-Південь (340), а також Амстердам (288). Лоукостери літають на курорти з багатьох аеропортів обох країн, тому граф містить багато пар «Німеччина $\to$ курорт $\to$ Велика Британія». Пошук у графі не враховує розкладу та авіакомпаній, тому такі маршрути -- можливі комбінації рейсів, а не квитки з пересадкою, які продають перевізники.

\subsection{Найвпливовіші аеропорти за PageRank}

П'ять найвпливовіших аеропортів за PageRank (\texttt{resetProbability = 0.15}): Атланта (ATL, 15,95), Чикаго О'Хара (ORD, 14,63), Стамбул Ататюрк (ISL, 14,62), Денвер (DEN, 14,54) та Даллас/Форт-Ворт (DFW, 14,29). Значення PageRank у GraphFrames нормовано так, що їх сума дорівнює кількості вершин (середнє -- 1), тобто ранг Атланти приблизно в 16 разів більший за середній.

\textbf{Чому саме вони.} PageRank аеропорту -- сума часток рангу аеропортів, з яких до нього є рейси, причому кожен аеропорт ділить свій ранг порівну між усіма своїми напрямками. Американські хаби (ATL, ORD, DEN, DFW) отримують ранг від сотень регіональних аеропортів США, для яких хаб -- один з небагатьох напрямків, тож вони віддають йому значну частку рангу. Крім того, хаби США густо пов'язані між собою і передають ранг один одному. Стамбул поєднує велику кількість напрямків (228 вхідних) з роллю транзитного вузла між Європою, Азією та Африкою.

\textbf{Порівняння зі ступеневою центральністю.} За кількістю напрямків лідирують FRA (239), CDG (237) і AMS (232), а за PageRank вони лише на 7--10 місцях (рис.~\ref{fig:pagerank},~\ref{fig:degree}). Європейські хаби мають багато напрямків, але значна частина їхніх сусідів -- інші великі хаби, які ділять свій ранг між сотнями напрямків, тож кожне ребро приносить мало рангу. Різниця показує, що PageRank оцінює не кількість, а «якість» зв'язків.

\textbf{Порівняння з найбільшими хабами світу.} Маршрути OpenFlights відповідають 2014~року, тому для порівняння взято рейтинг Airports Council International (ACI) за пасажиропотоком 2014~року~\cite{aci2014} (табл.~\ref{tab:aci}).

\begin{table}[!htb]
  \caption{Найбільші аеропорти світу за пасажиропотоком 2014~року (ACI) та їх місце за PageRank}\label{tab:aci}
  \begin{tabularx}{\textwidth}{|r|l|Y|r|r|}
    \hline
    Місце ACI & Код & Аеропорт & Пасажирів, млн & Місце за PageRank \\
    \hline
    1 & ATL & Атланта & 96,2 & 1 \\
    2 & PEK & Пекін Столичний & 86,1 & 9 \\
    3 & LHR & Лондон Хітроу & 73,4 & поза 10 \\
    4 & HND & Токіо Ханеда & 72,8 & поза 10 \\
    5 & LAX & Лос-Анджелес & 70,7 & поза 10 \\
    6 & DXB & Дубай & 70,5 & поза 10 \\
    7 & ORD & Чикаго О'Хара & 70,0 & 2 \\
    8 & CDG & Париж Шарль-де-Голль & 63,8 & 7 \\
    9 & DFW & Даллас/Форт-Ворт & 63,6 & 5 \\
    10 & HKG & Гонконг & 63,1 & поза 10 \\
    \hline
  \end{tabularx}
\end{table}

Збіг частковий. З п'ятірки PageRank до десятки ACI входять ATL (1), ORD (7) і DFW (9). Стамбул Ататюрк у 2014~році був 13-м за пасажиропотоком~\cite{aci2014-prelim}, Денвер до першої десятки не входив. Водночас Хітроу, Ханеда, Лос-Анджелес, Дубай і Гонконг -- 3--6 і 10 місця ACI -- не входять навіть до першої десятки за PageRank, а Хітроу -- і до десятки за кількістю напрямків (у нього менше ніж 187).

Розбіжність зумовлена тим, що саме вимірює PageRank на цьому графі. Граф містить одне ребро на пару аеропортів, без частоти рейсів, місткості літаків і кількості пасажирів, тому PageRank оцінює мережу напрямків, а не пасажиропотік. Щоденні рейси A380 між Лондоном і Дубаєм і рейс невеликого літака раз на тиждень мають однакову вагу. Тому вгору піднімаються аеропорти з густою внутрішньою мережею США, де хаб сполучений з сотнями регіональних аеропортів. Аеропорти з обмеженою кількістю слотів (LHR, HND) обслуговують величезну кількість пасажирів на меншій кількості напрямків -- частими рейсами великих літаків, а Ханеда у 2014~році обслуговувала переважно внутрішні рейси Японії. PageRank таких аеропортів нижчий, хоча за пасажиропотоком вони серед найбільших. Отже, PageRank на графі OpenFlights виявляє вузли, найважливіші для зв'язності мережі напрямків, і для оцінки хабів за пасажиропотоком потрібен граф з вагами ребер (частота рейсів або кількість місць).

\textbf{Вплив resetProbability.} Зі зростанням імовірності скидання значення PageRank вирівнюються: ранг лідера зменшується з 17,7 при 0,05 до 5,2 при 0,85, бо більша частина рангу розподіляється рівномірно між усіма вершинами. Склад лідерів також змінюється (рис.~\ref{fig:reset}). При малих значеннях (0,05) ранг поширюється по довгих ланцюжках, і до п'ятірки входять великі міжнародні хаби CDG та FRA -- оцінка глобальна. При великих значеннях (0,5--0,85) ранг вершини визначає переважно її найближче оточення, і лідерами стають DME, DEN, ATL -- аеропорти з великою кількістю регіональних напрямків, для яких вони є одним із небагатьох сполучень (Домодєдово для аеропортів Росії та СНД, Денвер для регіону Скелястих гір). Значення 0,15 -- компроміс, за якого враховується і транзитивний вплив, і локальна структура; при цьому ATL лідирує для значень 0,05--0,3.

\subsection{Кластери та ізольовані підграфи}

\textbf{Компоненти зв'язності.} Мережа майже повністю зв'язна: 3\,303 з 3\,329 аеропортів (99,2\,\%) утворюють одну слабко зв'язну компоненту. Решта 26 аеропортів розподілені на 6 ізольованих груп (табл.~\ref{tab:isolated}); одиночних ізольованих аеропортів немає, бо вершинами графа є лише аеропорти з маршрутами. Умову «не менше 5 аеропортів» задовольняють лише дві компоненти: основна мережа та група Нової Каледонії (10 аеропортів).

\begin{table}[H]
  \caption{Ізольовані підграфи}\label{tab:isolated}
  \begin{tabularx}{\textwidth}{|Y|r|Y|}
    \hline
    Регіон & Аеропортів & Аеропорти \\
    \hline
    Нова Каледонія & 10 & GEA (Нумеа), острови: BMY, ILP, KNQ, KOC, LIF та ін. \\
    Намібія & 4 & ERS (Віндгук), MPA, NDU, OND \\
    Алеутські острови (Аляска, США) & 4 & AKB, DUT, IKO, KQA \\
    Сієтл та острови Сан-Хуан (США) & 4 & BFI, CLM, ESD, FRD \\
    Американські Віргінські острови & 2 & SPB, SSB \\
    Невада -- Гранд-Каньйон (США) & 2 & BLD, GCW \\
    \hline
  \end{tabularx}
\end{table}

Ізольовані групи -- місцеві мережі, що мають лише внутрішні рейси: острівні сполучення (Нова Каледонія, Алеутські та Віргінські острови), регіональні рейси Намібії, туристичні рейси до Гранд-Каньйону. У Нумеа та Віндгуку внутрішні рейси виконуються з другого, «внутрішнього» аеропорту міста (GEA, ERS), а міжнародні -- з іншого, тому в графі ці мережі відірвані від основної, хоча пасажир може пересісти, переїхавши між аеропортами: міжнародний аеропорт Нумеа NOU належить основній компоненті (див. спільноту з NOU на рис.~\ref{fig:lpa-summary}). Частина ізольованості може бути й наслідком неповноти бази OpenFlights: граф містить лише маршрути, які є в ній на 2014~рік.

\textbf{Спільноти Label Propagation.} Оскільки основна компонента охоплює 99\,\% аеропортів, кластери всередині неї виділено алгоритмом Label Propagation: 71 спільнота з не менше ніж 5 аеропортів. Найбільші з них (табл.~\ref{tab:communities}) відповідають регіональним мережам: міжнародним хабам разом з частиною мережі США (796 аеропортів), Північній і Центральній Європі (351), Південній Європі з курортами Середземномор'я (336), Росії, СНД і Південній Азії (216), Східній Азії (159), Південно-Східній Азії та Океанії (157), внутрішній мережі США (102), регіональним аеропортам Китаю (85) та внутрішній мережі Бразилії (73).

\begin{table}[H]
  \caption{Найбільші спільноти Label Propagation}\label{tab:communities}
  \begin{tabularx}{\textwidth}{|r|r|l|Y|}
    \hline
    Аеропортів & Країн & Головна країна (частка) & Головні аеропорти \\
    \hline
    796 & 85 & США (0,43) & ISL, ATL, DXB, JFK, EWR \\
    351 & 50 & Норвегія (0,12) & FRA, CDG, AMS, MUC, LHR \\
    336 & 64 & Франція (0,10) & BCN, MAD, PMI, AGP, MXP \\
    216 & 33 & Індія (0,25) & DME, BKK, VKO, SVX, TAS \\
    159 & 20 & Японія (0,35) & PEK, PVG, CAN, HKG, ICN \\
    157 & 17 & Австралія (0,46) & SIN, KUL, SYD, MNL, CGK \\
    102 & 3 & США (0,49) & ORD, DFW, DEN, IAH, LAX \\
    85 & 3 & Китай (0,98) & WUX, TXN, CZX, XUZ, MIG \\
    73 & 1 & Бразилія (1,00) & BSB, CNF, SSA, CGH, CWB \\
    \hline
  \end{tabularx}
\end{table}

Спільноти мають чітку географічну основу: менші з них майже повністю належать одній країні (Китай 98\,\%, Бразилія 100\,\%, Аляска -- спільнота з BET, 100\,\% США), а великі об'єднують регіони навколо спільних хабів. Розбиття не збігається з межами країн: аеропорти США потрапили до трьох спільнот (глобальної з ATL і JFK, внутрішньої з ORD і DFW, аляскинської), а Європа поділена на «північну» мережу мережевих перевізників (FRA, AMS, LHR) і «південну» курортну мережу лоукостерів (BCN, PMI, AGP). Label Propagation виконано з фіксованою кількістю ітерацій (тут 5), без гарантії збіжності, а нічиї між мітками розв'язуються довільно, тому межі між великими спільнотами наближені; стабільним є саме поділ на регіональні мережі.

\FloatBarrier

% ============================================================================================
\anonsection{Висновки}

У роботі розгорнуто кластер Apache Spark~4.0.4 (master і два робочі вузли) у Docker і з бібліотекою GraphFrames проаналізовано мережу авіасполучень OpenFlights: 67\,663 маршрути, 12\,668 аеропортів і 6\,162 авіакомпанії.

Підготовка даних з'єднанням маршрутів з аеропортами дала координати для 67\,235 маршрутів. Перевірка найдовших відстаней виявила пошкоджений запис аеропорту 5613 (координати Каліфорнії для аеропорту в Замбії), маршрути через який виключено. Найбільшу сумарну довжину власних маршрутів має Ryanair (3,70~млн~км на 2\,484 маршрутах) -- за рахунок кількості маршрутів, а не їх довжини; найменшу -- TransHolding System (36~км на 2 маршрутах).

Граф з 3\,329 аеропортів і 37\,269 напрямків дозволив знайти маршрути з Німеччини до Великої Британії: мотиви знаходять 73 прямі маршрути, 9\,178 з однією пересадкою і 506\,498 з двома, тоді як BFS -- лише 73 найкоротші. BFS простіший у використанні (один виклик), але відповідає лише на питання про мінімальну кількість пересадок; мотиви потребують шаблону для кожної довжини та явних умов простого шляху, і їхній результат зростає комбінаторно. Між TXL та EDI прямого рейсу немає; найкоротший маршрут TXL $\to$ AMS $\to$ EDI має 1\,245~км, найдовший з не більше ніж трьома пересадками -- 37\,808~км.

Найвпливовіші за PageRank аеропорти -- ATL, ORD, ISL, DEN, DFW. Вони лише частково збігаються з найбільшими аеропортами світу за пасажиропотоком 2014~року: граф без ваг оцінює мережу напрямків, тому вгору піднімаються хаби густої внутрішньої мережі США, а аеропорти з обмеженою кількістю слотів і великим пасажиропотоком (LHR, HND) мають нижчий ранг. Зі зростанням \texttt{resetProbability} значення PageRank вирівнюються, а лідерами стають аеропорти з багатьма регіональними напрямками (DME, DEN).

Аналіз зв'язності показав, що 99,2\,\% аеропортів утворюють одну компоненту; 6 ізольованих груп (Нова Каледонія, Намібія, Алеутські острови та ін.) мають лише місцеві рейси. Label Propagation виділив 71 спільноту з не менше ніж 5 аеропортів, що відповідають регіональним мережам. Вебінтерфейс Spark підтвердив, що GraphFrames виконує графові алгоритми як запити DataFrame: мотив -- як ланцюжок broadcast-з'єднань з перенесеними оптимізатором фільтрами, компоненти зв'язності -- як ітеративні запити з checkpoint-ами, а PageRank -- через GraphX.

\printbibliography[heading=bibintoc, title={Список використаних джерел}]

% ============================================================================================
\clearpage
\anonsection{Додаток А. Програмний код}

Повний програмний код лабораторної роботи розміщено в репозиторії \url{https://github.com/nuinashco/iasa-bigdata}, директорія \texttt{labs/002-spark-graphframes}:
\begin{itemize}
  \item \texttt{draft.ipynb} -- ноутбук Jupyter з усіма кроками роботи та результатами виконання;
  \item \texttt{src/lab002/paths.py} -- пошук простих шляхів мотивами (додаток~Б);
  \item \texttt{src/lab002/spark.py} -- налаштування сесії Spark (драйвер у процесі ноутбука, пакет GraphFrames, checkpoint-директорія);
  \item \texttt{src/lab002/cluster.py} -- запуск кластера та виконання команд;
  \item \texttt{scripts/download\_data.py} -- завантаження даних OpenFlights (додаток~Б);
  \item \texttt{infra} -- конфігурація кластера (Docker Compose, образ Spark з Python~3.13).
\end{itemize}

Ноутбук відтворює роботу повністю: запуск «Run All» розгортає кластер, завантажує дані та виконує всі обчислення.

\anonsection{Додаток Б. Допоміжні модулі}

\setcounter{lstlisting}{0}
\renewcommand{\thelstlisting}{Б.\arabic{lstlisting}}
\lstinputlisting[language=Python, caption={Модуль paths.py: пошук простих шляхів мотивами}, label={lst:paths}]{../src/lab002/paths.py}

\lstinputlisting[language=Python, caption={Скрипт download\_data.py}, label={lst:download}]{../scripts/download_data.py}

\end{document}
